%==============================================================================
\section{Confusion Matrix (Ma trận nhầm lẫn)}
\label{sec:confusion-matrix}
%==============================================================================

\begin{dinhnghia}[title={Confusion Matrix}]
Cách \textbf{dễ nhất} đo hiệu năng bài toán phân loại khi đầu ra có hai hoặc nhiều lớp.
Bảng hai chiều ``Actual'' và ``Predicted'' với TP, TN, FP, FN.
\end{dinhnghia}

\tsimg{08_confusion_matrix/img01.jpg}

\begin{vidu}[title={Ví dụ spam email}]
Spam = \textbf{positive} (1); email hợp lệ = \textbf{negative} (0).
\end{vidu}

\begin{tinhchat}[title={True Positive (TP)}]
Thực tế và dự đoán đều 1; mô hình đúng lớp positive (ví dụ email spam đúng là spam).
\end{tinhchat}

\begin{tinhchat}[title={True Negative (TN)}]
Thực tế và dự đoán đều 0; đúng lớp negative.
\end{tinhchat}

\begin{tinhchat}[title={False Positive (FP)}]
Thực tế 0, dự đoán 1 --- \textbf{Type I error} (ví dụ email hợp lệ bị gán spam).
\end{tinhchat}

\begin{tinhchat}[title={False Negative (FN)}]
Thực tế 1, dự đoán 0 --- \textbf{Type II error} (spam bị gán hợp lệ).
\end{tinhchat}

\begin{phuongphap}[title={Phân loại đúng / sai}]
Đúng: TP và TN. Sai: FP và FN.
Từ confusion matrix tính accuracy, precision, recall, v.v.
\end{phuongphap}

\subsection{Ví dụ thực hành}

\begin{phuongphap}[title={10 email --- nhãn}]
Positive (spam) = 1; negative (not spam) = 0.

\textbf{Actual:} \quad 0\; 1\; 0\; 1\; 1\; 0\; 0\; 1\; 1\; 1

\textbf{Predicted:} \quad 0\; 1\; 0\; 1\; 0\; 1\; 0\; 0\; 1\; 1

\textbf{Kết quả:} \quad TN\; TP\; TN\; TP\; FN\; FP\; TN\; FN\; TP\; TP
\end{phuongphap}

\begin{phuongphap}[title={Đếm}]
TP = 4; FN = 2; FP = 1; TN = 3.

\begin{center}
\begin{tabular}{c|cc}
 & \textbf{Actual +} & \textbf{Actual $-$} \\
\hline
\textbf{Pred +} & 4 (TP) & 1 (FP) \\
\textbf{Pred $-$} & 2 (FN) & 3 (TN) \\
\end{tabular}
\end{center}

Trong 10 email: 7 đúng (TP+TN), 3 sai (FP+FN).
\end{phuongphap}

\subsection{Metric từ confusion matrix}

\begin{tinhchat}[title={Accuracy}]
\[
\mathrm{Accuracy} = \frac{TP+TN}{TP+FP+FN+TN} = \frac{7}{10} = 0.7
\]
\end{tinhchat}

\begin{tinhchat}[title={Precision}]
\[
\mathrm{Precision} = \frac{TP}{TP+FP} = \frac{4}{5} = 0.8
\]
\end{tinhchat}

\begin{tinhchat}[title={Recall (Sensitivity)}]
\[
\mathrm{Recall} = \frac{TP}{TP+FN} = \frac{4}{6} \approx 0.667
\]
\end{tinhchat}

\begin{tinhchat}[title={Specificity}]
\[
\mathrm{Specificity} = \frac{TN}{TN+FP} = \frac{3}{4} = 0.75
\]
\end{tinhchat}

\begin{tinhchat}[title={F1 Score}]
\[
F1 = 2 \times \frac{Precision \times Recall}{Precision + Recall} \approx 0.727
\]
\end{tinhchat}

\begin{tinhchat}[title={Type I Error Rate}]
\[
\frac{FP}{FP+TN} = \frac{1}{4} = 0.25
\]
\end{tinhchat}

\begin{tinhchat}[title={Type II Error Rate}]
\[
\frac{FN}{FN+TP} = \frac{2}{6} \approx 0.333
\]
\end{tinhchat}

\subsection{Triển khai Python}

\begin{luuy}[title={Thứ tự ma trận sklearn}]
\texttt{confusion\_matrix(y\_actual, y\_pred)} trả về:
\begin{center}
\begin{tabular}{c|cc}
 & Pred 0 & Pred 1 \\
\hline
Actual 0 & TN & FP \\
Actual 1 & FN & TP \\
\end{tabular}
\end{center}
\end{luuy}

\begin{vidu}[title={confusion\_matrix}]
\begin{lstlisting}
from sklearn.metrics import confusion_matrix

y_actual = [0, 1, 0, 1, 1, 0, 0, 1, 1, 1]
y_pred   = [0, 1, 0, 1, 0, 1, 0, 0, 1, 1]
cm = confusion_matrix(y_actual, y_pred)
print(cm)
\end{lstlisting}

\textbf{Output:}
\begin{verbatim}
[[3 1]
 [2 4]]
\end{verbatim}
Khớp bảng ví dụ trên (TN=3, FP=1, FN=2, TP=4).
\end{vidu}

\begin{vidu}[title={Heatmap với seaborn}]
\begin{lstlisting}
import seaborn as sns
sns.heatmap(cm, annot=True, cmap='summer')
\end{lstlisting}
\tsimg{08_confusion_matrix/img02.jpg}

Trục $x$: dự đoán; trục $y$: thực tế; màu thể hiện số mẫu mỗi ô.
\end{vidu}
